\documentclass[11pt]{article}
\usepackage[T1]{fontenc}
\usepackage{lmodern,amsmath,amssymb,amsthm}
\usepackage[margin=1in]{geometry}
\usepackage[hidelinks]{hyperref}
\newtheorem{theorem}{Theorem}
\newtheorem{lemma}[theorem]{Lemma}
\newtheorem{corollary}[theorem]{Corollary}
\newcommand{\Q}{\mathbb Q}
\newcommand{\Z}{\mathbb Z}
\newcommand{\F}{\mathbb F}
\newcommand{\Sha}{\mathrm{Sha}}
\newcommand{\Sel}{\operatorname{Sel}}
\newcommand{\ord}{\operatorname{ord}}
\title{A Kurihara Witness and the Leading Values of All Full-Support Twists}
\author{\small Project correspondence: \href{mailto:admin@cosmosgeek.org}{admin@cosmosgeek.org}}
\date{Public working paper for specialist review}
\begin{document}\maketitle
\begin{abstract}
For a weight-two newform attached to an elliptic curve, a single nonzero
mod-$p$ modular-symbol moment determines the exact cyclotomic valuation
of every central twist value whose local characters have the same
$p$-power order and the prescribed full support. The proof uses an
inclusion--exclusion identity before evaluating the norm relations.
There is no bound on the number of prime factors. We obtain an exact
valuation $hd$ for the relative leading value in a cyclic field of degree
$p^h$ with $d$ primes in the conductor. The Selmer application additionally
requires Kim's hypotheses, $d$ independent rational points, and cyclic
local $p$-groups of sufficient order at the ramified primes.
\end{abstract}

\section{A fixed-conductor congruence}
Let $E/\Q$ have conductor $N$, with normalized newform
$f=\sum_{n\ge1}a_n e^{2\pi inz}$ and real period $\Omega^+$.
Use the plus symbol
\[
 c(x)=\frac{\operatorname{Re}\left(2\pi\int_0^\infty f(x+iy)\,dy\right)}{\Omega^+}.
\]
The period convention in Kim \cite[\S1.4]{Kim} is the integral of the
absolute N\'eron differential over $E(\mathbb R)$.
Fix an odd prime $p$. Let $n=\ell_1\cdots\ell_d$, $d\ge1$, be
squarefree with $(n,Np)=1$. Assume $c(a/m)\in\Z_{(p)}$ for
$m\mid n$ and $(a,m)=1$, including $c(0)$.
For $h\ge1$ assume
\[
 \ell_i\equiv1\pmod{p^h},\qquad a_{\ell_i}\equiv2\pmod p .
 \tag{1}
\]
Choose a primitive root $\eta_i$ at each $\ell_i$ and let $j_i(a)$
be the discrete logarithm of $a$ to this base. Put
\[
 \delta_n=\sum_{a\in(\Z/n)^\times}c(a/n)
         \prod_{i=1}^d j_i(a)\pmod p.
\]
The condition $\delta_n\ne0$ is independent of the choices of
primitive roots.

\begin{theorem}\label{main}
Let $1\le b\le h$ and $\zeta$ be a primitive $p^b$th root of unity.
Let $\chi=\prod_i\chi_i$ be a Dirichlet character with
$\chi_i(\eta_i)=\zeta^{u_i}$, $p\nmid u_i$ for every $i$.
Put $\tau(\chi)=\sum_{a\bmod n}\chi(a)e^{2\pi ia/n}$,
$\pi=\zeta-1$, and normalize $v_\pi(\pi)=1$. Then
\[
 T(\chi):=\frac{\tau(\chi)L(E,\overline\chi,1)}{\Omega^+}
 \in\pi^d\Z_{(p)}[\zeta],\qquad
 \frac{T(\chi)}{\pi^d}\equiv\delta_n\prod_i u_i\pmod\pi .
 \tag{2}
\]
In particular $\delta_n\ne0$ implies $v_\pi(T(\chi))=d$ for every such
character.
\end{theorem}

\begin{proof}
For $m\mid n$ and a character $\psi$ modulo $m$, write
\[
 A_m(\psi)=\sum_{a\in(\Z/m)^\times} c(a/m)\psi(a),\qquad A_1(1)=c(0).
\]
For primitive $\psi$ of odd order the Gauss-sum identity gives
$A_m(\psi)=T(\psi)$. These characters are even, so the minus modular
symbols cancel in pairs. Fourier expansion and
$\sum_a\psi(a)e^{2\pi ika/m}=\overline\psi(k)\tau(\psi)$ give the
stated convention, by Abel limiting at each of the finitely many cusps.
For $\ell\mid n/m$ and a character inflated from $m$,
\[
 A_{m\ell}(\psi)=
 (a_\ell-\psi(\ell)-\psi(\ell)^{-1})A_m(\psi).
 \tag{3}
\]
Indeed the weight-two Hecke relation on the primitive of $f$ gives
$\sum_{j\bmod\ell} c((x+j)/\ell)=a_\ell c(x)-c(\ell x)$.
For $x=a/m$ the unique term with numerator divisible by $\ell$
is $c(\ell^{-1}a/m)$, with the inverse read modulo $m$.
Removing it and summing proves (3), also for $m=1$.

The factor in (3) belongs to $(\pi^2)$ by (1) and
$\xi+\xi^{-1}-2=(\xi-1)^2/\xi$ for $\xi=\psi(\ell)$.
Here $v_\pi(p)=\varphi(p^b)\ge2$. For a subset $I$ of the
$d$ primes, let $m_I$ be their product and $\chi_I$ the product
of the corresponding local characters. Consider
\[
 \sum_{a\in(\Z/n)^\times}c(a/n)\prod_{i=1}^d(\chi_i(a)-1)
   =\sum_{I\subseteq\{1,\dots,d\}}(-1)^{d-|I|}A_n(\chi_I).
 \tag{4}
\]
The left side divided by $\pi^d$ reduces to
$\delta_n\prod_i u_i$. We prove (2) by induction on $d$, using
only the lower bound $v_\pi(A_{m_I}(\chi_I))\ge|I|$ for proper
subsets. For the empty set this is integrality of $c(0)$.
Iteration of (3) then gives
$v_\pi(A_n(\chi_I))\ge |I|+2(d-|I|)\ge d+1$
for every proper subset. The full subset in (4) is $T(\chi)$.
\end{proof}

\section{The relative leading term}
Let $\chi$ have exact local order $p^h$ at every $\ell_i$ and let $F$
be the cyclic degree-$p^h$ field cut out by $\chi$. There are
$\varphi(p^h)^{d-1}$ distinct such fields, since there are
$\varphi(p^h)^d$ characters and each kernel has $\varphi(p^h)$ faithful characters.
\begin{corollary}\label{relative}
If $\delta_n\ne0$, both Mordell--Weil rank and complex analytic rank
are unchanged from $\Q$ to $F$. Moreover
\[
 B_F=(\Omega^+)^{1-p^h}
 \lim_{s\to1}\frac{L(E/F,s)}{L(E/\Q,s)}
 \in\Q^\times,\qquad v_p(B_F)=hd .
 \tag{5}
\]
\end{corollary}
\begin{proof}
For each $1\le b\le h$, the $\varphi(p^b)$ order-$p^b$ characters of
$\operatorname{Gal}(F/\Q)$ have full conductor $n$ and equal local
orders. All their central values are nonzero by (2). Artin
factorization proves the equality of vanishing orders. For $E$ without
complex multiplication, Kato's theorem gives
$(E(F)\otimes\mathbb C)^\psi=0$ for each nontrivial character
component with nonzero central value
\cite[\S14]{Kato}, \cite[III, Thm.~5.4(ii)]{Rubin}.
For $E$ with complex multiplication the corresponding conclusion
comes from elliptic units \cite[III, Remark~5.7]{Rubin}.
The $T(\psi)$ in each exact-order orbit are algebraic conjugates.
The sum of their $v_p$ is $\varphi(p^b)\,d/\varphi(p^b)=d$.
The Gauss sums over that orbit multiply to $n^{\varphi(p^b)/2}$,
since $\tau(\psi)\tau(\overline\psi)=n$ for these even characters.
Their product is a rational $p$-unit. Multiplication over the $h$
orbits gives (5).
\end{proof}
If $a=\ord_{s=1}L(E,s)$, simultaneous comparison at $1+u/p^h$
upstairs and $1+u$ downstairs instead multiplies (5) by $p^{-ha}$.
The resulting valuation is $h(d-a)$. The norm relation (3) proves
(5) at a common argument and imposes no integrality on the rescaled
ratio.

\section{Rank certification and the local norm conditions}
For the rest of this section take $p\ge5$ and $\Omega^+$ in Kim's
N\'eron normalization. Assume that the Galois representation
$G_{\Q}\to\mathrm{GL}_2(\Z_p)$ on $T_pE$ is surjective, the Manin
constant is prime to $p$, $E(\Q_p)[p]=0$, and every Tamagawa number is
prime to $p$. Suppose that (1) holds, $\delta_n\ne0$, and $E(\Q)$
contains $d$ independent points. Then Kim's Theorems 1.8 and 1.10
\cite{Kim} imply
\[
 \operatorname{rank}E(\Q)=d,\qquad \Sha(E/\Q)[p^\infty]=0 .
 \tag{6}
\]
Indeed Theorem 1.8 identifies the Selmer corank with the minimum number
of primes of a nonzero Kurihara number, at most $d$ and at least the
rational rank. The same theorem's Fitting formula gives the unit
$d$th Fitting ideal, since the minimum valuation at weight $d$ and
the overall minimum are both zero. The dual Selmer group thus has rank
$d$ and is free on $d$ generators. The Kummer exact sequence gives
a surjection from this free module onto the dual of
$E(\Q)\otimes\Q_p/\Z_p$, also free of rank $d$. Its kernel
$\Sha(E/\Q)[p^\infty]^\vee$ is zero. Further, the minimum weight of a mod-$p$ unit is
$d$, since a unit also gives a nonzero Kurihara number. Kim's
Theorem 1.10 gives the semilocal isomorphism at this witness.

\begin{corollary}
In addition to all hypotheses of this section, suppose that for each
$\ell\mid n$ the group
$E(\F_\ell)[p^\infty]$ is cyclic of order divisible by $p^h$.
Then $\Sha(E/F)[p^\infty]=0$.
\end{corollary}
\begin{proof}
Write $G=\operatorname{Gal}(F/\Q)$. The irreducible residual image
has no fixed vectors over $F$. By normality its fixed subspace
is either zero or all, and the full residual image cannot factor
through a $p$-group. By Corollary~\ref{relative}, the ranks are
equal. Thus $E(F)\otimes\Z_p=E(\Q)\otimes\Z_p$. A difference of a
Galois conjugate of a point and that point is prime-to-$p$ torsion,
and multiplication by the prime-to-$p$ torsion exponent descends it.
In particular $H^1(G,E(F))=0$ and the $p$-part of
$H^2(G,E(F))$ is identified by cyclic cohomology with
$E(\Q)/p^hE(\Q)$.

At $\ell\mid n$ the extension is totally ramified of degree $p^h$.
Its norm cokernel is $E(\F_\ell)/p^hE(\F_\ell)$. Reduction of the
norm acts by $p^h$, and the reduction kernel is uniquely
$p$-divisible. At an unramified finite place the norm is
surjective on the identity component. On its finite residue group
this follows from Lang's theorem and cyclic cohomology, and on each
successive formal-group quotient from surjectivity of the residue-field
trace. Successive approximation lifts to the formal group. The remaining
component group at the base has order prime to $p$ by hypothesis.
Real places split. By local Tate duality the restriction kernels on
local Weil--Ch\^atelet groups thus have no $p$-part outside $n$.

The resulting diagonal map
\[
 E(\Q)/p^hE(\Q)\longrightarrow
 \bigoplus_{\ell\mid n}E(\F_\ell)/p^hE(\F_\ell)
\]
has source and target free of rank $d$ over $\Z/p^h$. Its reduction modulo $p$
is an isomorphism by (6) and Kim's semilocal isomorphism, so it
is an isomorphism over $\Z/p^h$.
If $\beta\in\Sha(E/F)[p]^G$, its transgression in $H^2(G,E(F))$
vanishes locally. Injectivity of the global-to-local norm map
makes that transgression zero. Thus $\beta$ descends to
$c\in H^1(\Q,E)[p]$ by inflation--restriction and $H^1(G,E(F))=0$.
The localizations of $c$ vanish outside $n$. At the primes of $n$
they pair trivially in sum with every $Q\in E(\Q)$ by global
reciprocity for the Tate pairing on $E$ and its dual (identified by the principal
polarization). The surjective diagonal map and local
duality force each localization to vanish. Hence $c=0$ by (6).
A nonzero finite $\F_p[G]$-module has nonzero invariants when $G$
is a $p$-group; applying this to $\Sha(E/F)[p]$ proves the claim.
\end{proof}

\section{A published witness and its twists}
Kim gives $p=5$, $n=41\cdot61$ and $\delta_n\ne0$ for the curve
$389\mathrm{a}1$, with equation $y^2+y=x^3+x^2-2x$.
The same source records rational rank two, trivial $5$-primary
Tate--Shafarevich group over $\Q$, and the semilocal localization
isomorphism \cite[\S8, example (1)]{Kim}.
For $q=41,61$ the point-count formula
\[
 \#E(\F_q)=1+q+\sum_{x\bmod q}
       \left(\frac{1+4(x^3+x^2-2x)}{q}\right)
\]
gives $45$ and $70$, respectively. In particular $a_{41}=-3$,
$a_{61}=-8$, and each local $5$-primary group has order five.
Theorem~\ref{main} gives $v_{\zeta_5-1}(T(\chi))=2$ for all
$4^2=16$ full-support characters. Their kernels define four cyclic
quintic fields. Equation (5) gives $v_5(B_F)=2$ for each field.


The horizontal nonvanishing results of Kriz and Nordentoft
\cite[Cor.~5.20]{KN} give estimates across conductors. The
calculation here evaluates the leading residue at a given unit
witness, for every equal-order full-support character and its lower
order powers. Birch--Stevens interpolation, the Hecke norm
relation, and the results of Kim, Kato and Rubin provide the inputs.
This is a project working paper prepared with ChatGPT and Codex.
Independent literature-priority review remains open.
\begin{thebibliography}{9}
\bibitem{Kim} C.-H. Kim, \emph{The structure of Selmer groups and
the Iwasawa main conjecture for elliptic curves}, Theorems 1.8 and 1.10.
\href{https://preprint.press.jhu.edu/ajm/sites/default/files/AJM-kim.pdf}{Paper}.
\bibitem{KN} D. Kriz and A. C. Nordentoft, \emph{Horizontal $p$-adic
$L$-functions}, Corollary~5.20.
\href{https://arxiv.org/abs/2310.20678}{Paper}.
\bibitem{Kato} K. Kato, \emph{$p$-adic Hodge theory and values of zeta
functions of modular forms}, Ast\'erisque 295 (2004), 117--290.
\href{https://www.numdam.org/item/AST_2004__295__117_0/}{Paper}.
\bibitem{Rubin} K. Rubin, \emph{Euler Systems}, Chapter~III,
Theorem~5.4 and Remark~5.7. \href{https://swc-math.github.io/notes/files/99RubinES.pdf}{Lecture notes}.
\end{thebibliography}
\end{document}
